\section{Results III: The Mechanism and Causal Test}

\subsection{Predictions Ledger}

\begin{table}[h]
\centering
\begin{tabular}{p{0.16\linewidth}p{0.42\linewidth}p{0.28\linewidth}}
\toprule
Test & Prediction & Measured \\
\midrule
Exp1 & Supervision does not give intermediate volume control & Dead units in every baseline seed \\
Exp1 & Partial volume control is worse than none & Redundancy increases from 20.5 to 187.0 \\
Exp3 & Optimizer signature disappears under composition & SGD spans 83.8--96.2; Adam helps \\
Exp4-P1 & CE dominates EM at small $\lambda$ & 30--70x, cosine near zero \\
Exp4-P2 & LSE curvature stays below $1/2$ & Bound held; CE path 13--100x larger \\
Exp4-P3 & Stop-gradient restores conditioning & LR spread falls from 10.9 to 1.3 points \\
\bottomrule
\end{tabular}
\caption{Prediction ledger for the supervised study.}
\label{tab:predictions-ledger}
\end{table}

\subsection{Experiment 4 Design}

The causal experiment compares four arms: joint training with
$\lambda=0.001$, joint training with $\lambda=0.03$, joint training with
$\lambda=1$, and stop-gradient training where the supervised head reads
detached intermediate distances. Feature quality is measured by a fixed-protocol
linear probe trained on frozen distances.

\subsection{Causal Restoration of Conditioning}

Probe accuracy spread across a $1000\times$ SGD learning-rate range falls from
10.9 points in the CE-dominated joint arm to 1.3 points under stop-gradient.
The $\lambda=1$ joint arm is indistinguishable from stop-gradient, while
$\lambda=0.03$ lies between the two regimes. Supervision is not the problem; the
composition corridor is.

\subsection{Gradient Competition}

At $\lambda=0.001$, the CE gradient reaching the intermediate distances
dominates the weighted auxiliary EM gradient by $30$--$70\times$. Their cosine
similarity is approximately zero. The EM signal is drowned, not opposed.

\subsection{Curvature Dichotomy}

The LSE Hessian with respect to $d$ remains bounded by $1/2$ at every measured
point. Under pure EM training it saturates the bound, reaching
$0.4993$--$0.4997$. The CE-path Hessian is roughly $13$--$100\times$ larger,
non-monotone, and consistent with $W_2$ spectral scaling. Distance from
$1/2$ therefore acts as a diagnostic of mixture sharpness.

\subsection{Lambda-Resolution Curve}

The graded claim should be supported by a denser sweep over $\lambda$. The
planned central figure plots conditioning spread and probe accuracy against the
measured CE/EM gradient ratio, not against $\lambda$ itself.

\subsection{Basin Analysis}

Linear mode connectivity between EM-dominated and CE-dominated checkpoints will
formalize the fine-tuning claim. Hidden units must be permutation-aligned before
interpolation, for example with Hungarian matching.

\subsection{Adam Arm Under Stop-Gradient}

The Adam arm completes the Paper 2 optimizer signature. The prediction is that
Adam provides little advantage under stop-gradient and a clear advantage in the
joint CE-dominated arm.

\subsection{Stochastic-Map Arm}

Replacing $W_2$ with a non-negative row-normalized map tests the one
simplex-compatible preservation class not yet covered experimentally. A positive
result would provide the controlled bridge to attention's row-stochastic value
path.
